\section{Global deformation theorems}
\label{sec:global-theorems}

The physical matrix of Section~\ref{sec:compact-gauging} has an independent
geometric characterization. We formulate it in terms of the topology of the
singularity links and Milnor fibers. This section also distinguishes the exact
first-order statement from the all-orders condition on an analytic smoothing
arc.

\subsection{Links, Milnor fibers and branch lattices}
\label{subsec:link-lattices}

Let $\Sigma_p$ be the link of the isolated singularity $(X,p)$. For a
smoothing component $\sigma_p$, let $F_{p,\sigma_p}$ be a Milnor fiber and
define the vanishing subspace
\begin{equation}
 V_{p,\sigma_p}
 =\ker\!\left[
 H_2(\Sigma_p;\QQ)\longrightarrow H_2(F_{p,\sigma_p};\QQ)
 \right].
 \label{eq:milnor-kernel}
\end{equation}
The spaces used in this paper are
\begin{equation}
 \begin{array}{c|c|c}
 (X,p)&H_2(\Sigma_p;\QQ)&V_{p,\sigma_p}\\
 \hline
 \text{node}&\QQ&\QQ\\
 C(S_7)&K_{S_7}^{\perp}\otimes\QQ&A_1\otimes\QQ\\
 C(S_6),\ \cH_{\mathfrak{sl}_2}&K_{S_6}^{\perp}\otimes\QQ&A_1\otimes\QQ\\
 C(S_6),\ \cH_{\mathfrak{sl}_3}&K_{S_6}^{\perp}\otimes\QQ&A_2\otimes\QQ.
 \end{array}
 \label{eq:local-milnor-table}
\end{equation}
For a del Pezzo cone, the link is the unit-circle bundle of $K_{S_p}$ and
the Gysin sequence gives
\begin{equation}
 H_2(\Sigma_p;\QQ)\cong K_{S_p}^{\perp}\subset H_2(S_p;\QQ),
 \label{eq:link-gysin}
\end{equation}
where the superscript $\perp$ refers to the intersection pairing. The root
subspaces in \eqref{eq:local-milnor-table} are the cycles that die
in the corresponding Milnor fiber (Proposition~2.1 and Section~2.4 of
\cite{Wuebben:2026obstruction}, where the three sweeping Fano threefolds
are identified).

Choose a branch profile $\sigma$ and set
\begin{equation}
 \Kbranch
 =\ker\!\left[
   \bigoplus_pH_2(\Sigma_p;\QQ)\longrightarrow H_2(\wideX;\QQ)
 \right]
 \cap\bigoplus_pV_{p,\sigma_p}.
 \label{eq:topological-kernel}
\end{equation}
The map sends a nodal generator to its exceptional curve class and a del
Pezzo link class to its pushforward from the del Pezzo exceptional divisor. Thus
\eqref{eq:topological-kernel} is the rational kernel of
\eqref{eq:geometric-charge-map}.

\begin{lemma}[local deformation--link comparison]
\label{lem:local-comparison}
Let $\widehat X_p\to(X,p)$ be the chosen crepant local resolution, with
exceptional set $E_p$, and put
\begin{equation}
 K_p^{\mathrm{cr}}=
 \ker\!\left[
 H^2_{E_p}(\widehat X_p;\Omega^2_{\widehat X_p})
 \longrightarrow H^2(\widehat X_p;\Omega^2_{\widehat X_p})
 \right].
 \label{eq:local-crepant-kernel}
\end{equation}
For a node and for the cones over $S_7$ and $S_6$, contraction with the
crepant volume form and the boundary map give natural isomorphisms
\begin{equation}
 T^1_{X,p}\cong K_p^{\mathrm{cr}}
 \cong H_2(\Sigma_p;\CC).
 \label{eq:local-hodge-identification}
\end{equation}
These isomorphisms commute with the maps to
$H^2(\wideX;\Omega^2_{\wideX})$ and $H_2(\wideX;\CC)$ induced by the local
inclusions. They carry the tangent cone of a smoothing component to
$V_{p,\sigma_p}\otimes\CC$.
\end{lemma}

For completeness, the branch assertion is finite. The circle-bundle Gysin
sequence identifies the del Pezzo link with $K_{S_p}^{\perp}$. On the
pentagon, the reflection fixing the unique smoothing parameter acts, after
the determinant twist from contraction with the volume form, on the unique
$A_1$ root line. On the hexagon, a sixfold rotation has a $(-1)$ eigenspace
and a Coxeter eigenspace; these are the $A_1$ line and $A_2$ plane. The local
cohomology sequence identifies these eigenspaces with
\eqref{eq:local-crepant-kernel}. This proves the branch statement and the
compatibility with the Gysin maps; the full derivation is Lemmas~8.6
and~8.7 of \cite{Wuebben:2026smoothing}, which rest on Schlessinger's
comparison \cite{Schlessinger:1971} in the depth-three form of
\cite{Friedman:2022}.

\subsection{A class of toric anticanonical hypersurfaces}
\label{subsec:admissible-class}

Let $N\cong\ZZ^4$, let $M=\operatorname{Hom}(N,\ZZ)$, and write
$\langle\ ,\ \rangle$ for the natural pairing. For a four-dimensional
reflexive polytope $\Delta\subset N_{\mathbb R}$, set
\begin{equation}
 \Delta^\circ=\{u\in M_{\mathbb R}:\langle u,v\rangle\geq-1
 \text{ for every }v\in\Delta\}.
 \label{eq:polar-convention}
\end{equation}
Let $\Sigma_\Delta$ be the fan of cones over the faces of $\Delta$ and
$\mathbb P_\Delta=\operatorname{TV}(\Sigma_\Delta)$ its Gorenstein Fano
toric fourfold. A hypersurface $X\subset\mathbb P_\Delta$ is
$\Delta$-regular if its Laurent equation and every face restriction are
nondegenerate in the sense of Batyrev \cite{Batyrev:1994}.

If $G\subset\Delta$ is a two-face, its dual face $G^*\subset\Delta^\circ$
is an edge. A $\Delta$-regular anticanonical hypersurface meets the torus
orbit associated with $G$ in $\ell(G^*)$ points, where $\ell$ is lattice
length; the transverse germ at each point is the affine cone over the toric
surface defined by $G$. We call $G$ \emph{singular} when it is not a
unimodular triangle. The unit square, the unit-edge pentagon with one interior
lattice point, and the analogous hexagon give, respectively, a node and the
cones over $S_7$ and $S_6$.

The explicit finite theorem is proved for the following controlled class. We
call $\Delta$ admissible if:
\begin{enumerate}
 \item every two-face is a unimodular triangle, a unit square, or a
 pentagon or hexagon with one interior lattice point;
 \item every edge has lattice length one; and
 \item the dual edge of every singular two-face has lattice length one.
\end{enumerate}
For a $\Delta$-regular anticanonical hypersurface, the singular locus is then
finite and consists exactly of nodes and cones over $S_7$ or $S_6$. The
ambient divisor matrix detects the full rational topological kernel in
\eqref{eq:topological-kernel}; equivalently, the correction term in
Batyrev's formula for $h^{1,1}$ \cite{Batyrev:1994} vanishes on this class
(Proposition~8.3 of \cite{Wuebben:2026obstruction}).

The abstract local-to-global statement below uses the following additional
hypotheses:
\begin{equation}
 \omega_X\cong\cO_X,
 \qquad h^1(X,\cO_X)=0,
 \qquad
 \wideX\text{ satisfies the }\partial\bar\partial\text{-lemma}.
 \label{eq:global-hypotheses}
\end{equation}
The third holds automatically for the projective crepant resolution fixed
in Section~\ref{subsec:m2-charges}; all three hold in the toric
hypersurface examples of Section~\ref{sec:examples}.

\subsection{The first-order image theorem}
\label{subsec:first-order-theorem}

Let $\TangentDef$ denote the global first-order deformation space of $X$. We
view $V_{p,\sigma_p}\otimes\CC$ as a subspace of $T^1_{X,p}$ through
\eqref{eq:local-hodge-identification}.

\begin{theorem}[first-order local-to-global image]
\label{thm:first-order-image}
Let $X$ arise from an admissible polytope and satisfy
\eqref{eq:global-hypotheses}. For every branch profile $\sigma$,
\begin{equation}
 \operatorname{Im}\!\left[
   \TangentDef\longrightarrow\bigoplus_pT^1_{X,p}
 \right]
 \cap\bigoplus_p\bigl(V_{p,\sigma_p}\otimes\CC\bigr)
 =\Kbranch\otimes\CC.
 \label{eq:first-order-image-theorem}
\end{equation}
Equivalently, in the divisor and root bases of
Section~\ref{sec:compact-gauging}, the right-hand side is
$\ker\Qcompact$.
\end{theorem}

\begin{proof}
For a normal Gorenstein threefold with finite
non-local-complete-intersection (non-lci) locus, the cotangent
complex and Kähler-differential deformation groups agree through degree two
(Lemma~8.5 of \cite{Wuebben:2026smoothing}).
The local-to-global Ext sequence therefore begins
\begin{equation}
 0\longrightarrow H^1(X;\Theta_X)\longrightarrow\TangentDef
 \xrightarrow{\ \operatorname{loc}\ }\bigoplus_pT^1_{X,p}
 \xrightarrow{\ \mathrm{ob}\ }H^2(X;\Theta_X).
 \label{eq:local-to-global-ext}
\end{equation}
We identify the kernel of $\mathrm{ob}$ on each branch.

Let $E=\operatorname{Exc}(\wideX\to X)$ and
$X_{\mathrm{reg}}=\wideX\setminus E$. Excision and the local-cohomology
sequence give the exact segment
\begin{equation}
 H^1(X_{\mathrm{reg}};\Omega^2_{X_{\mathrm{reg}}})\longrightarrow
 \bigoplus_p H^2_{E_p}(\widehat X_p;\Omega^2_{\widehat X_p})
 \longrightarrow H^2(\wideX;\Omega^2_{\wideX}).
 \label{eq:global-support-sequence}
\end{equation}
By Lemma~\ref{lem:local-comparison}, each branch space
$V_{p,\sigma_p}\otimes\CC$ is a subspace of $K_p^{\mathrm{cr}}$ in the
middle term. Moreover, the middle map in
\eqref{eq:global-support-sequence} fits into the commutative square
\begin{equation}
\begin{array}{ccc}
 \displaystyle\bigoplus_pH_2(\Sigma_p;\CC)&\longrightarrow&H_2(\wideX;\CC)\\
 \downarrow\wr&&\downarrow\operatorname{PD}\\
 \displaystyle\bigoplus_pK_p^{\mathrm{cr}}&\longrightarrow&H^4(\wideX;\CC),
\end{array}
 \label{eq:gysin-comparison-square}
\end{equation}
where the lower arrow is the last map of
\eqref{eq:global-support-sequence} followed by the de Rham map. Thus a link
class in the topological kernel \eqref{eq:topological-kernel} has zero de
Rham image. The $\partial\bar\partial$-lemma makes
$H^2(\wideX;\Omega^2_{\wideX})\to H^4(\wideX;\CC)$ injective, so exactness of
\eqref{eq:global-support-sequence} lifts that class to
$H^1(X_{\mathrm{reg}};\Omega^2_{X_{\mathrm{reg}}})$.

Contraction with a trivialization of $\omega_X$ identifies
$H^1(X_{\mathrm{reg}};\Omega^2_{X_{\mathrm{reg}}})$ with
$H^1(X_{\mathrm{reg}};\Theta_{X_{\mathrm{reg}}})$. Schlessinger's
depth-three comparison \cite{Schlessinger:1971}, in the form of
\cite{Friedman:2022}, identifies the latter with $\TangentDef$, and
naturality identifies
the resulting localization map with $\operatorname{loc}$ in
\eqref{eq:local-to-global-ext}. This proves that every element of
$\Kbranch\otimes\CC$ lies in the left-hand side of
\eqref{eq:first-order-image-theorem}. Conversely, any global deformation
class maps to zero under $\mathrm{ob}$; the same commutative square sends its
link class to zero in $H_2(\wideX;\CC)$. Lemma~\ref{lem:local-comparison}
then identifies its branch restriction with an element of $\Kbranch$. The
two inclusions prove the equality. This argument is the kernel comparison of
Theorem~D (Theorem~8.8) and Section~8.4 of \cite{Wuebben:2026smoothing},
printed here to make the maps and obstruction group explicit.
\end{proof}

Combining Theorem~\ref{thm:first-order-image} with
Proposition~\ref{prop:physical-kernel} yields the central identification
\begin{equation}
 \operatorname{Im}\!\left(
 \cM_{H,\sigma}^{\mathrm{rigid}}\to\cB_{\sigma}
 \right)
 =\operatorname{Im}\!\left(
 \TangentDef\to\bigoplus_pT^1_{X,p}
 \right)\cap\cB_{\sigma}.
 \label{eq:physical-geometric-identification}
\end{equation}
This is the sense in which the geometric kernel is physically derived as a
five-dimensional charge and gauging constraint.

\subsection{All-orders necessity}
\label{subsec:all-orders}

A first-order statement does not by itself constrain an analytic arc whose
local parameter begins at higher order. The necessary condition for a genuine
smoothing is obtained instead from vanishing periods. For a node, a
distinguished vanishing three-sphere has period proportional to the local
deformation parameter. For the divisorial contractions used here, the Milnor
fiber contains a distinguished three-cycle whose period is nonzero on the
relevant component. Mayer--Vietoris boundary maps relate these cycles to
classes in \eqref{eq:topological-kernel}. The resulting argument is insensitive
to the order of contact with the local discriminant. The required local
period statement and its simultaneous full-support consequence are
Theorem~6.1 and Remark~6.2, respectively, of
\cite{Wuebben:2026obstruction}.

\begin{theorem}[all-orders necessary relation]
\label{thm:all-orders-necessity}
Let $X$ be a compact complex threefold with
$\omega_X\cong\cO_X$ and $h^1(X,\cO_X)=0$, whose singularities are nodes and
anticanonical cones over $S_7$ or $S_6$. Fix a crepant resolution
$\wideX\to X$, small over the nodes and with the corresponding del Pezzo
surface over each cone point. Let
$\mathcal X\to\mathbb D$ be a flat proper analytic family over a disc, with
central fiber $\mathcal X_0\cong X$ and smooth fibers $\mathcal X_s$ for all
sufficiently small $s\ne0$, and let $\sigma$ be the branch profile it
induces: at each $E_3$ point, the component of the miniversal base through
which the family smooths the germ. Then the projection
\begin{equation}
 \Kbranch\longrightarrow V_{q,\sigma_q}
 \label{eq:nonzero-projection}
\end{equation}
is nonzero for every node $q$, every $E_2$ point $q$, and every $E_3$ point
$q$ at which $\sigma_q=\mathfrak{sl}_2$.
\end{theorem}

\begin{remark}
\label{rem:full-support}
Since $\QQ$ is infinite, a generic rational linear combination avoids the
finite union of the coordinate hyperplanes in question, so $\Kbranch$
contains a single class whose component is nonzero simultaneously at every
germ named in Theorem~\ref{thm:all-orders-necessity}. The statement applies
to the actual analytic arc and not merely to its first derivative.
\end{remark}

Theorem~\ref{thm:all-orders-necessity} gives no conclusion at an $E_3$
point smoothed through $\cH_{\mathfrak{sl}_3}$, although it still constrains
every other germ of the same family, with $\Kbranch$ computed on the profile
that family induces. That component has a two-dimensional Milnor kernel. The
analogous proof requires two independent period pairings, not only a nonzero
scalar period, and remains open. The first-order
Theorem~\ref{thm:first-order-image} does include this plane.

\subsection{Integrability boundary}
\label{subsec:integrability}

Theorem~\ref{thm:first-order-image} implies that every branch-restricted class
in the kernel lifts to a global first-order deformation. It does not imply
that this deformation integrates. The del Pezzo cone germs are not local
complete intersections \cite{Gross:1997}, and the standard unobstructedness
argument for nodal
Calabi--Yau threefolds does not apply uniformly on the non-$\QQ$-factorial
locus relevant here.

Three distinct assertions should therefore remain separate:
\begin{enumerate}
 \item $b\in\ker\Qcompact$ solves the projected rigid moment-map equations;
 \item $b$ is the localization of a global first-order complex deformation;
 \item $b$ is tangent to an actual finite extremal transition.
\end{enumerate}
The first two are proved for every profile in the admissible class. The third
is proved for purely nodal transitions and for the explicit $X_9$ family in
Section~\ref{subsec:x9}; it is open in general. This is the remaining step in
a necessary-and-sufficient mixed analogue of Friedman's conifold-transition
criterion \cite{Friedman:1986}.
